\bookchapter{第六章：守护你的数字生活——安全与隐私}

\chapterintro{手机安全=隐私安全+财产安全。在这个数据即资产的时代，保护手机就是保护你自己。}

\booksection{一、锁屏：第一道物理防线}

锁屏是你手机的第一道防线，也是最容易被忽视的。

\booksubsection{安全等级排序：}

\begin{enumerate}
  \item \textbf{密码/PIN码}：安全度最高，是最终保障（忘记指纹/脸时靠它）。
  \item \textbf{指纹识别}：安全与便利的完美平衡。
  \item \textbf{人脸识别}：最方便，但安全性取决于技术（3D结构光安全，2D照片易骗过）。
  \item \textbf{图案解锁}：方便但容易留下痕迹被猜到。
\end{enumerate}

\booksubsection{设置强锁屏：}

\begin{itemize}
  \item 务必设置一个复杂密码（不要用生日、123456）。
  \item 开启生物识别。
  \item 锁屏时间设成"立即"或最短。
\end{itemize}

\booksection{二、“查找我的手机”：终极防盗工具}

\textbf{必须开启！}它能：

\begin{itemize}
  \item \textbf{定位}：手机丢哪了，在地图上找。
  \item \textbf{播放声音}：让手机响铃，适合在家找。
  \item \textbf{丢失模式}：远程锁死手机并显示联系号码。
  \item \textbf{远程擦除}：最后一招，清除所有数据（但擦除后无法再定位）。
\end{itemize}

去 [设置] 里，苹果找“查找我的iPhone”，安卓找“查找我的设备”或品牌类似服务，确保它是绿色开启状态。

\booksection{三、系统更新：打补丁，堵漏洞}

系统更新首要目的是修复安全漏洞，就像给房子补墙洞。所以，看到系统更新提示，尽量第一时间升级。

当然，大版本更新（比如iOS 17到iOS 18）可能会有兼容性问题，可以观望几天再升。

\booksection{四、识别与防范电信诈骗}

在AI技术迅猛发展的今天，传统的“眼见为实”认知模式已完全失效。诈骗分子正利用深度伪造（Deepfake）等AI技术，实现了从“广撒网”式诈骗到“精准围猎”式犯罪的全面升级。

\booksubsection{AI 时代最常见的几种诈骗套路}

\textbf{(1) AI 换脸拟声诈骗}\par

不法分子通过非法搜集目标对象在社交媒体发布的高清照片或仅需几秒钟的语音样本，利用AI深度伪造技术克隆出足以乱真的“亲友视频”或“语音”。常见剧本包括：\textbf{突发车祸、紧急借款、遭遇绑架}等。此类诈骗利用亲情、友情等情感纽带，\textbf{迷惑性极高}。

\textbf{(2) 冒充权威 / 客服诈骗}\par

AI技术能够精准模拟官方客服、公检法机关人员的声音与语调，甚至可以合成高度逼真的虚假转账截图、电子公文、通缉令等文件，通过\textbf{制造恐慌心理或权威压力}，诱导受害者转账或开启屏幕共享功能，从而窃取账户信息。

\textbf{(3) AI 投资理财诈骗}\par

通过AI算法伪造“稳赚不赔”的虚假收益截图、虚拟量化选股软件界面或“大师级”投资分析报告，以高额回报为诱饵，售卖高价投资服务或诱导受害者将资金转入其控制的交易账户，\textbf{最终卷款跑路}。

\booksubsection{破局之法——“照妖镜”防护思路}

AI可以克隆声音和面容，但也不是无懈可击，有几个破绽可以利用：

\begin{itemize}
  \item \textbf{动作验证法}：进行视频通话时，主动要求对方完成预设的连贯动作指令，例如"用手在脸前快速挥过"、"侧脸展示耳朵轮廓"或"连续快速眨眼5次"。目前的AI换脸技术在处理复杂动态遮挡、快速运动场景和光影变化时，\textbf{极易出现画面卡顿、面部扭曲或边缘穿模}等破绽。
  \item \textbf{暗号核验法}：建立并使用预设的私密核验问题，例如"我们上次一起聚餐时你点的什么菜"、"我送你的第一个礼物是什么"等只有双方知晓的细节。若对方出现明显迟疑或答非所问，应\textbf{立即终止通话并启动安全预案}。
  \item \textbf{二次核实法}：严格遵守公安部网安局发布的"\textbf{四不一核实}"原则——不轻信陌生信息、不随意转账汇款、不泄露个人敏感信息、不传播可疑内容。涉及任何钱财往来时，务必通过\textbf{已知的原号码回拨或当面确认}。进行转账操作时，建议选择"24小时到账"模式，为潜在诈骗识别和资金挽回预留充足时间窗口。
\end{itemize}

\booksubsection{传统常见套路（仍要警惕）：}

\begin{itemize}
  \item \textbf{中奖/刷单返利}：利诱，先给你点甜头再骗大钱。
  \item \textbf{网恋诈骗}：杀猪盘，培养感情后骗钱。
  \item \textbf{冒充熟人}：盗号后借钱（AI 时代更难分辨）。
\end{itemize}

\booksubsection{手机自带防护：}

\begin{itemize}
  \item 骚扰电话拦截。
  \item 验证码防窃取。
  \item 短信链接安全检测。
\end{itemize}

\booksubsection{防骗口诀：}

\textbf{不轻信、不透露、不点击、不转账、有疑问先核实！}遇到“急事”先冷静5分钟，几乎所有诈骗都会在这一刻露出马脚。

\booksection{五、抵御“开盒”与社工攻击：斩断隐私泄露链条}

“开盒”（社会工程学攻击）的核心原理在于\textbf{系统性搜集、分析并拼凑个人数字碎片，构建完整的个人画像}。防范的关键在于为个人隐私划定明确红线，建立多层次防护体系。

\booksubsection{照片脱敏处理}

基础 EXIF 清理方法见第二章“图片与文件系统管理”。在社交媒体分享照片前，还需进一步\textbf{彻底删除照片的元数据（EXIF信息）}，包括拍摄时间、设备型号、GPS位置坐标等敏感数据。对照片中包含的\textbf{门牌号码、车牌信息、车票条形码、身份证号}等关键信息进行严格打码或模糊处理，避免无意中泄露个人行踪轨迹。

\begin{itemize}
  \item iOS：[设置] > [隐私与安全性] > [相机] 中可关闭"精确位置"；分享时选"仅当前照片"可剥离元数据。
  \item 安卓：部分相册App（如Google相册）分享时会提示是否保留位置信息。
  \item 通用工具：ExifTool、Metapho 等可查看/编辑 EXIF。
\end{itemize}

\booksubsection{应用权限收敛}

基础权限管理方法详见第四章“权限管理”。在“防开盒”场景下，还需进一步定期审查并关闭电商类、工具类App的\textbf{麦克风、通讯录、位置信息}等非必要权限，防止App在后台“偷听”日常对话或非法收集个人关系网络；对长期闲置但仍占用高权限的应用程序应及时卸载，从源头减少隐私泄露风险点。

\begin{itemize}
  \item 看到"通讯录权限"申请时，先问一句：这App真有必要知道我所有联系人吗？
  \item 三个月没用过的App，果断卸载。权限是越用越紧的，不是越用越松。
\end{itemize}

\booksubsection{智能设备安全设防}

智能摄像头、智能门锁、智能音箱等IoT设备在初次使用时必须\textbf{立即修改默认登录密码}，设置包含大小写字母、数字和特殊符号的高强度独立密码；不常用的远程访问功能应保持关闭状态，定期检查设备固件更新并及时修补安全漏洞，\textbf{防止黑客入侵并窥探家庭隐私空间}。

\begin{tcbwarn}{真实案例警示}
2022年某地警方通报，黑客通过破解某品牌家用摄像头的默认密码（很多用户从未修改），窃取了上千户家庭的卧室画面并在网上兜售。\textbf{不改默认密码 = 把家大门敞开}。

\end{tcbwarn}

\booksection{六、隐私保护：你的数据是你的财产}

手机里的照片、聊天记录、位置轨迹，本质上都是你的“数字资产”。一旦泄露，轻则广告骚扰，重则被人肉、诈骗甚至敲诈。

\booksubsection{App 权限的“最小化原则”：}

\begin{itemize}
  \item \textbf{位置}：打车/地图用"始终允许"；其他 App 用"仅使用时允许"。
  \item \textbf{通讯录}：除非必要（如微信加好友）才给，否则一律禁止。
  \item \textbf{相册 / 全部照片}：iOS 14+ 起可授权"选定的照片"，安卓 13+ 同样支持精细授权。
  \item \textbf{麦克风 / 摄像头}：平时关掉，需要时再开启。怀疑被偷录时，滑出控制中心看是否有"小绿点/小橙点"。
  \item \textbf{后台跟踪}：iOS 14.5+ 起强制要求 App 使用 IDFA 前必须询问；安卓可关闭"广告 ID"。
\end{itemize}

\booksubsection{隐私管理实战：}

\begin{itemize}
  \item \textbf{隐私报告（iOS）}：[设置] > [隐私与安全性] > [App 隐私报告]，查看哪些 App 在偷偷访问你的位置、相机、麦克风。
  \item \textbf{权限使用记录（安卓）}：[设置] > [隐私] > [权限使用记录]，看谁在什么时候调用了敏感权限。
  \item \textbf{限制广告跟踪}：在系统设置里关闭"广告 ID"或"限制广告追踪"，能减少"被精准广告骚扰"。
\end{itemize}

\booksection{七、丢失 / 被盗后的“急救”清单}

手机丢了不要慌，按以下顺序处理：

\begin{enumerate}
  \item \textbf{立即冻结 SIM 卡}：联系运营商挂失，避免被用于接收验证码。
  \item \textbf{远程锁定}：通过云服务（iCloud / 品牌云）开启"丢失模式"。
  \item \textbf{修改重要密码}：登录过手机的所有重要账号（支付、银行、邮箱）立即改密码。
  \item \textbf{报案}：到派出所报案，记录 IMEI 号码。手机找回时也便于证明身份。
  \item \textbf{通知亲友}：防止骗子以你名义借钱。
  \item \textbf{（最后）远程擦除}：确认找回无望时再操作，擦除后定位功能失效。
\end{enumerate}

\booksection{八、SIM 卡与 eSIM 安全}

\begin{itemize}
  \item \textbf{设置 SIM 卡 PIN}：开启后，每次重启手机或换卡都需要输入 PIN 码，防止"拔卡换手机"绕过锁屏。
  \item \textbf{eSIM 优势}：没有实体卡可被拔出，安全性更高；同时支持多号切换，旅行时非常方便。
  \item \textbf{二次换卡风险}：警惕"运营商二次换卡"骗局——骗子冒用你的身份补办 SIM 卡，从而接管你的手机号。设置"补卡需机主本人持身份证到营业厅办理"是最有效的防线。
\end{itemize}

\begin{tcbwarn}{千万别做的事}
\begin{itemize}
  \item 用生日 / 123456 / qwerty 等弱密码；
  \item 把验证码发给任何人（哪怕是"客服"）；
  \item 点击短信里的不明链接，特别是 00 开头的境外号码；
  \item 在公共 Wi-Fi 下登录网银或支付；
  \item 给陌生人开启"屏幕共享"或"远程控制"。
\end{itemize}

\end{tcbwarn}

\begin{tcbpractice}{试着启用并模拟"远程锁定"}
\begin{enumerate}
  \item 确保你的"查找我的手机"功能已开启。
  \item 在电脑浏览器或另一部手机上，登录对应的云服务账号（iCloud.com 或品牌云服务网站）。
  \item 找到"查找设备"功能，尝试对你自己的手机发起一次"播放声音"或模拟"锁定"指令，熟悉整个操作流程。
\end{enumerate}

\end{tcbpractice}
